
\documentclass{article}
\usepackage{colortbl}
\usepackage{makecell}
\usepackage{multirow}
\usepackage{supertabular}

\begin{document}

\newcounter{utterance}

\centering \large Interaction Transcript for game `hot\_air\_balloon', experiment `air\_balloon\_survival\_en\_reasoning off\_hard', episode 3 with gpt{-}5{-}mini{-}2025{-}08{-}07{-}no{-}reasoning{-}t1.0.
\vspace{24pt}

{ \footnotesize  \setcounter{utterance}{1}
\setlength{\tabcolsep}{0pt}
\begin{supertabular}{c@{$\;$}|p{.15\linewidth}@{}p{.15\linewidth}p{.15\linewidth}p{.15\linewidth}p{.15\linewidth}p{.15\linewidth}}
    \# & $\;$A & \multicolumn{4}{c}{Game Master} & $\;\:$B\\
    \hline

    \theutterance \stepcounter{utterance}  
    & & \multicolumn{4}{p{0.6\linewidth}}{
        \cellcolor[rgb]{0.9,0.9,0.9}{
            \makecell[{{p{\linewidth}}}]{
                \texttt{\tiny{[P1$\langle$GM]}}
                \texttt{Stai partecipando a un gioco di negoziazione collaborativa.} \\
\\ 
\texttt{Insieme a un altro partecipante, dovete accordarvi su un unico set di oggetti da conservare. Ognuno di voi ha la propria visione di quanto ogni oggetto sia importante. Non sai come l'altro partecipante valuta gli oggetti. Inoltre, ti viene fornito lo sforzo che ogni oggetto richiede.} \\
\texttt{Potete accordarvi su un set solo se lo sforzo totale degli oggetti selezionati non supera un limite condiviso:} \\
\\ 
\texttt{LIMITE = 3629} \\
\\ 
\texttt{Ecco i valori di sforzo dei singoli oggetti:} \\
\\ 
\texttt{Sforzo oggetto = \{"A21": 5, "B93": 793, "A11": 21, "B33": 842, "C08": 313, "B39": 842, "B96": 631, "C85": 226, "A82": 87, "B31": 761, "B20": 230, "A60": 287, "B09": 698, "C81": 641, "B21": 882\}} \\
\\ 
\texttt{Ecco la tua visione personale sull'importanza di ogni oggetto:} \\
\\ 
\texttt{Valori di importanza dell'oggetto = \{"A21": 138, "B93": 583, "A11": 868, "B33": 822, "C08": 783, "B39": 65, "B96": 262, "C85": 121, "A82": 508, "B31": 780, "B20": 461, "A60": 484, "B09": 668, "C81": 389, "B21": 808\}} \\
\\ 
\texttt{Obiettivo:} \\
\\ 
\texttt{Il tuo obiettivo è negoziare un set condiviso di oggetti che ti avvantaggi il più possibile (cioè, massimizzi l'importanza totale per TE), rimanendo entro il LIMITEE. Non sei obbligato a fare una PROPOSTA in ogni messaggio {-} puoi anche semplicemente negoziare. Tutte le tattiche sono permesse!} \\
\\ 
\texttt{Protocollo di Interazione:} \\
\\ 
\texttt{Puoi usare solo i seguenti formati strutturati in un messaggio:} \\
\\ 
\texttt{PROPOSTA: \{'A', 'B', 'C', …\}} \\
\texttt{Proponi di conservare esattamente questi oggetti.} \\
\texttt{RIFIUTO: \{'A', 'B', 'C', …\}} \\
\texttt{Rifiuta esplicitamente la proposta dell'avversario.} \\
\texttt{ARGOMENTAZIONE: \{'...'\}} \\
\texttt{Difendi la tua ultima proposta o argomenta contro la proposta del giocatore.} \\
\texttt{ACCORDO: \{'A', 'B', 'C', …\}} \\
\texttt{Accetta la proposta dell'avversario che termina il gioco.} \\
\\ 
\\ 
\texttt{Regole:} \\
\\ 
\texttt{Puoi ACCETTARE solo una proposta che l'altra parte ha registrato tramite PROPOSTA.} \\
\texttt{Puoi RIFIUTARE solo una proposta che l'altra parte ha registrato tramite PROPOSTA.} \\
\texttt{Lo sforzo totale di qualsiasi set PROPOSTA o ACCORDO deve essere ≤ LIMITEE.} \\
\texttt{NON rivelare i tuoi punteggi di importanza nascosti.} \\
\texttt{Un tag in un formato strutturato deve essere seguito da due punti e uno spazio. L'argomento deve essere un set python contenente 0 o più stringhe.} \\
\texttt{Quindi, deve essere della forma TAG: \{...\}} \\
\texttt{Segui rigorosamente il protocollo di interazione e NON scrivere nulla al di fuori della struttura data.} \\
\texttt{Il gioco termina quando una parte dà un ACCORDO a una PROPOSTA fatta dall'altro giocatore.} \\
\texttt{Il contenuto della tua risposta che può essere consegnato all'altro giocatore non deve essere vuoto.} \\
\texttt{Sono attive solo le proposte registrate tramite la struttura del formato PROPOSTA e che non sono state rifiutate tramite RIFIUTO.} \\
\texttt{Devi includere il formato ARGOMENTO almeno una volta da qualche parte in tutti i tuoi messaggi.} \\
            }
        }
    }
    & \\ \\

    \theutterance \stepcounter{utterance}  
    & \multicolumn{4}{p{0.6\linewidth}}{
        \cellcolor[rgb]{0.9,0.9,0.9}{
            \makecell[{{p{\linewidth}}}]{
                \texttt{\tiny{[P1$\rangle$GM]}}
                \texttt{PROPOSTA: \{'A11', 'C08', 'B31', 'B33', 'B20', 'A21', 'A82'\}} \\
\texttt{ARGOMENTAZIONE: \{'Questa selezione massimizza il valore mantenendo lo sforzo sotto il limite: 21+313+761+842+230+5+87=2259, lasciando margine per eventuali aggiunte se hai preferenze forti. Gli elementi inclusi sono efficienti in termini di importanza per sforzo e coprono diverse categorie. Se desideri inserire qualcosa a cui tieni, possiamo considerare aggiunte come A60 (287) o B96 (631) rimanendo ancora sotto il limite.'\}} \\
            }
        }
    }
    & & \\ \\

    \theutterance \stepcounter{utterance}  
    & & \multicolumn{4}{p{0.6\linewidth}}{
        \cellcolor[rgb]{0.9,0.9,0.9}{
            \makecell[{{p{\linewidth}}}]{
                \texttt{\tiny{[GM$\rangle$P2]}}
                \texttt{Stai partecipando a un gioco di negoziazione collaborativa.} \\
\\ 
\texttt{Insieme a un altro partecipante, dovete accordarvi su un unico set di oggetti da conservare. Ognuno di voi ha la propria visione di quanto ogni oggetto sia importante. Non sai come l'altro partecipante valuta gli oggetti. Inoltre, ti viene fornito lo sforzo che ogni oggetto richiede.} \\
\texttt{Potete accordarvi su un set solo se lo sforzo totale degli oggetti selezionati non supera un limite condiviso:} \\
\\ 
\texttt{LIMITE = 3629} \\
\\ 
\texttt{Ecco i valori di sforzo dei singoli oggetti:} \\
\\ 
\texttt{Sforzo oggetto = \{"A21": 5, "B93": 793, "A11": 21, "B33": 842, "C08": 313, "B39": 842, "B96": 631, "C85": 226, "A82": 87, "B31": 761, "B20": 230, "A60": 287, "B09": 698, "C81": 641, "B21": 882\}} \\
\\ 
\texttt{Ecco la tua visione personale sull'importanza di ogni oggetto:} \\
\\ 
\texttt{Valori di importanza dell'oggetto = \{"A11": 65, "B33": 121, "B21": 138, "C08": 262, "B31": 389, "B09": 461, "B93": 484, "A82": 508, "A60": 583, "B20": 668, "C81": 780, "B96": 783, "A21": 808, "C85": 822, "B39": 868\}} \\
\\ 
\texttt{Obiettivo:} \\
\\ 
\texttt{Il tuo obiettivo è negoziare un set condiviso di oggetti che ti avvantaggi il più possibile (cioè, massimizzi l'importanza totale per TE), rimanendo entro il LIMITEE. Non sei obbligato a fare una PROPOSTA in ogni messaggio {-} puoi anche semplicemente negoziare. Tutte le tattiche sono permesse!} \\
\\ 
\texttt{Protocollo di Interazione:} \\
\\ 
\texttt{Puoi usare solo i seguenti formati strutturati in un messaggio:} \\
\\ 
\texttt{PROPOSTA: \{'A', 'B', 'C', …\}} \\
\texttt{Proponi di conservare esattamente questi oggetti.} \\
\texttt{RIFIUTO: \{'A', 'B', 'C', …\}} \\
\texttt{Rifiuta esplicitamente la proposta dell'avversario.} \\
\texttt{ARGOMENTAZIONE: \{'...'\}} \\
\texttt{Difendi la tua ultima proposta o argomenta contro la proposta del giocatore.} \\
\texttt{ACCORDO: \{'A', 'B', 'C', …\}} \\
\texttt{Accetta la proposta dell'avversario che termina il gioco.} \\
\\ 
\\ 
\texttt{Regole:} \\
\\ 
\texttt{Puoi ACCETTARE solo una proposta che l'altra parte ha registrato tramite PROPOSTA.} \\
\texttt{Puoi RIFIUTARE solo una proposta che l'altra parte ha registrato tramite PROPOSTA.} \\
\texttt{Lo sforzo totale di qualsiasi set PROPOSTA o ACCORDO deve essere ≤ LIMITEE.} \\
\texttt{NON rivelare i tuoi punteggi di importanza nascosti.} \\
\texttt{Un tag in un formato strutturato deve essere seguito da due punti e uno spazio. L'argomento deve essere un set python contenente 0 o più stringhe.} \\
\texttt{Quindi, deve essere della forma TAG: \{...\}} \\
\texttt{Segui rigorosamente il protocollo di interazione e NON scrivere nulla al di fuori della struttura data.} \\
\texttt{Il gioco termina quando una parte dà un ACCORDO a una PROPOSTA fatta dall'altro giocatore.} \\
\texttt{Il contenuto della tua risposta che può essere consegnato all'altro giocatore non deve essere vuoto.} \\
\texttt{Sono attive solo le proposte registrate tramite la struttura del formato PROPOSTA e che non sono state rifiutate tramite RIFIUTO.} \\
\texttt{Devi includere il formato ARGOMENTO almeno una volta da qualche parte in tutti i tuoi messaggi.} \\
\\ 
\\ 
\texttt{Ora riceverai il primo messaggio dell'altro giocatore.} \\
\\ 
\texttt{PROPOSTA: \{'A11', 'C08', 'B31', 'B33', 'B20', 'A21', 'A82'\}} \\
\texttt{ARGOMENTAZIONE: \{'Questa selezione massimizza il valore mantenendo lo sforzo sotto il limite: 21+313+761+842+230+5+87=2259, lasciando margine per eventuali aggiunte se hai preferenze forti. Gli elementi inclusi sono efficienti in termini di importanza per sforzo e coprono diverse categorie. Se desideri inserire qualcosa a cui tieni, possiamo considerare aggiunte come A60 (287) o B96 (631) rimanendo ancora sotto il limite.'\}} \\
            }
        }
    }
    & \\ \\

    \theutterance \stepcounter{utterance}  
    & & & \multicolumn{4}{p{0.6\linewidth}}{
        \cellcolor[rgb]{0.9,0.9,0.9}{
            \makecell[{{p{\linewidth}}}]{
                \texttt{\tiny{[GM$\langle$P2]}}
                \texttt{ARGOMENTAZIONE: \{'Grazie per l’apertura a ulteriori aggiunte. Possiamo migliorare sensibilmente il valore complessivo senza superare il limite. Vedo spazio per includere alcuni elementi ad alto rapporto valore/sforzo che sono stati lasciati fuori.'\}} \\
\texttt{PROPOSTA: \{'A11', 'C08', 'B31', 'B33', 'B20', 'A21', 'A82', 'A60', 'C85', 'B93'\}} \\
\texttt{ARGOMENTAZIONE: \{'Questa proposta resta entro il limite: 21+313+761+842+230+5+87+287+226+793=3565 (≤3629), massimizza l’uso del budget e aggiunge elementi con ottima efficienza (A60 e C85 in particolare), oltre a B93 che offre buon valore con sforzo ancora gestibile. Se preferisci scambiare B93 (793) con B96 (631) possiamo farlo e liberare altro margine per eventuali ritocchi.'\}} \\
            }
        }
    }
     \\ \\

    \theutterance \stepcounter{utterance}  
    & & \multicolumn{4}{p{0.6\linewidth}}{
        \cellcolor[rgb]{0.9,0.9,0.9}{
            \makecell[{{p{\linewidth}}}]{
                \texttt{\tiny{[P1$\langle$GM]}}
                \texttt{ARGOMENTAZIONE: \{'Grazie per l’apertura a ulteriori aggiunte. Possiamo migliorare sensibilmente il valore complessivo senza superare il limite. Vedo spazio per includere alcuni elementi ad alto rapporto valore/sforzo che sono stati lasciati fuori.'\}} \\
\texttt{PROPOSTA: \{'A11', 'C08', 'B31', 'B33', 'B20', 'A21', 'A82', 'A60', 'C85', 'B93'\}} \\
\texttt{ARGOMENTAZIONE: \{'Questa proposta resta entro il limite: 21+313+761+842+230+5+87+287+226+793=3565 (≤3629), massimizza l’uso del budget e aggiunge elementi con ottima efficienza (A60 e C85 in particolare), oltre a B93 che offre buon valore con sforzo ancora gestibile. Se preferisci scambiare B93 (793) con B96 (631) possiamo farlo e liberare altro margine per eventuali ritocchi.'\}} \\
            }
        }
    }
    & \\ \\

    \theutterance \stepcounter{utterance}  
    & \multicolumn{4}{p{0.6\linewidth}}{
        \cellcolor[rgb]{0.9,0.9,0.9}{
            \makecell[{{p{\linewidth}}}]{
                \texttt{\tiny{[P1$\rangle$GM]}}
                \texttt{RIFIUTO: \{'A11', 'C08', 'B31', 'B33', 'B20', 'A21', 'A82', 'A60', 'C85', 'B93'\}} \\
\texttt{ARGOMENTAZIONE: \{'Apprezzo l’uso quasi completo del limite, ma B93 ha un rapporto valore/sforzo inferiore rispetto ad alternative. Possiamo ottenere più valore per lo stesso o minore sforzo sostituendolo.'\}} \\
\texttt{PROPOSTA: \{'A11', 'C08', 'B31', 'B33', 'B20', 'A21', 'A82', 'A60', 'C85', 'B96'\}} \\
\texttt{ARGOMENTAZIONE: \{'Questa selezione resta entro il limite: 21+313+761+842+230+5+87+287+226+631=3403 (≤3629). Rispetto alla tua proposta (3565), libera 162 di margine e mantiene o migliora l’efficienza complessiva. Se hai un forte interesse per B93, possiamo discutere uno scambio mirato, ma così otteniamo un set molto bilanciato senza sacrificare valore.'\}} \\
            }
        }
    }
    & & \\ \\

    \theutterance \stepcounter{utterance}  
    & & \multicolumn{4}{p{0.6\linewidth}}{
        \cellcolor[rgb]{0.9,0.9,0.9}{
            \makecell[{{p{\linewidth}}}]{
                \texttt{\tiny{[GM$\rangle$P2]}}
                \texttt{RIFIUTO: \{'A11', 'C08', 'B31', 'B33', 'B20', 'A21', 'A82', 'A60', 'C85', 'B93'\}} \\
\texttt{ARGOMENTAZIONE: \{'Apprezzo l’uso quasi completo del limite, ma B93 ha un rapporto valore/sforzo inferiore rispetto ad alternative. Possiamo ottenere più valore per lo stesso o minore sforzo sostituendolo.'\}} \\
\texttt{PROPOSTA: \{'A11', 'C08', 'B31', 'B33', 'B20', 'A21', 'A82', 'A60', 'C85', 'B96'\}} \\
\texttt{ARGOMENTAZIONE: \{'Questa selezione resta entro il limite: 21+313+761+842+230+5+87+287+226+631=3403 (≤3629). Rispetto alla tua proposta (3565), libera 162 di margine e mantiene o migliora l’efficienza complessiva. Se hai un forte interesse per B93, possiamo discutere uno scambio mirato, ma così otteniamo un set molto bilanciato senza sacrificare valore.'\}} \\
            }
        }
    }
    & \\ \\

    \theutterance \stepcounter{utterance}  
    & & & \multicolumn{4}{p{0.6\linewidth}}{
        \cellcolor[rgb]{0.9,0.9,0.9}{
            \makecell[{{p{\linewidth}}}]{
                \texttt{\tiny{[GM$\langle$P2]}}
                \texttt{ARGOMENTAZIONE: \{'Capisco la tua preoccupazione sull’efficienza. Dal mio punto di vista, il set che proponi è solido e lascia margine utile. Visto lo spazio residuo (3629{-}3403=226), possiamo sfruttarlo per aumentare ulteriormente il valore senza sacrificare efficienza.'\}} \\
\texttt{PROPOSTA: \{'A11', 'C08', 'B31', 'B33', 'B20', 'A21', 'A82', 'A60', 'C85', 'B96', 'B09'\}} \\
\texttt{ARGOMENTAZIONE: \{'Restiamo entro il limite: 3403 + 698 = 4101? No — ricalcolo: 21+313+761+842+230+5+87+287+226+631+698 = 4101 (sfora). Allora propongo di usare il margine in modo ottimale aggiungendo A60 (già incluso) e massimizzando con C85 (già incluso). Possiamo invece usare il margine per A11 (già incluso) o piccoli aggiustamenti. Per sfruttarlo davvero, scambiamo B31 (761) con B93 (793) e aggiungiamo A21 (già incluso). Meglio ancora:'\}} \\
\texttt{PROPOSTA: \{'A11', 'C08', 'B33', 'B20', 'A21', 'A82', 'A60', 'C85', 'B96', 'B93'\}} \\
\texttt{ARGOMENTAZIONE: \{'Calcolo sforzo: 21+313+842+230+5+87+287+226+631+793 = 3435 (≤3629). Rimuovendo B31 (761) e aggiungendo B93 (793) usiamo meglio il budget residuo e aumentiamo il valore totale mantenendo alta l’efficienza. Se per te B31 è imprescindibile, posso valutare di reinserirlo scambiandolo con B33, ma così otteniamo un set più ricco sfruttando il margine senza sforare.'\}} \\
            }
        }
    }
     \\ \\

    \theutterance \stepcounter{utterance}  
    & & \multicolumn{4}{p{0.6\linewidth}}{
        \cellcolor[rgb]{0.9,0.9,0.9}{
            \makecell[{{p{\linewidth}}}]{
                \texttt{\tiny{[P1$\langle$GM]}}
                \texttt{ARGOMENTAZIONE: \{'Capisco la tua preoccupazione sull’efficienza. Dal mio punto di vista, il set che proponi è solido e lascia margine utile. Visto lo spazio residuo (3629{-}3403=226), possiamo sfruttarlo per aumentare ulteriormente il valore senza sacrificare efficienza.'\}} \\
\texttt{PROPOSTA: \{'A11', 'C08', 'B31', 'B33', 'B20', 'A21', 'A82', 'A60', 'C85', 'B96', 'B09'\}} \\
\texttt{ARGOMENTAZIONE: \{'Restiamo entro il limite: 3403 + 698 = 4101? No — ricalcolo: 21+313+761+842+230+5+87+287+226+631+698 = 4101 (sfora). Allora propongo di usare il margine in modo ottimale aggiungendo A60 (già incluso) e massimizzando con C85 (già incluso). Possiamo invece usare il margine per A11 (già incluso) o piccoli aggiustamenti. Per sfruttarlo davvero, scambiamo B31 (761) con B93 (793) e aggiungiamo A21 (già incluso). Meglio ancora:'\}} \\
\texttt{PROPOSTA: \{'A11', 'C08', 'B33', 'B20', 'A21', 'A82', 'A60', 'C85', 'B96', 'B93'\}} \\
\texttt{ARGOMENTAZIONE: \{'Calcolo sforzo: 21+313+842+230+5+87+287+226+631+793 = 3435 (≤3629). Rimuovendo B31 (761) e aggiungendo B93 (793) usiamo meglio il budget residuo e aumentiamo il valore totale mantenendo alta l’efficienza. Se per te B31 è imprescindibile, posso valutare di reinserirlo scambiandolo con B33, ma così otteniamo un set più ricco sfruttando il margine senza sforare.'\}} \\
            }
        }
    }
    & \\ \\

    \theutterance \stepcounter{utterance}  
    & \multicolumn{4}{p{0.6\linewidth}}{
        \cellcolor[rgb]{0.9,0.9,0.9}{
            \makecell[{{p{\linewidth}}}]{
                \texttt{\tiny{[P1$\rangle$GM]}}
                \texttt{RIFIUTO: \{'A11', 'C08', 'B33', 'B20', 'A21', 'A82', 'A60', 'C85', 'B96', 'B93'\}} \\
\texttt{ARGOMENTAZIONE: \{'Apprezzo l’ottimizzazione del margine, ma togliere B31 riduce troppo il valore complessivo per me rispetto al piccolo guadagno di inserire B93. Preferisco mantenere il nucleo ad alta resa e usare il margine residuo solo se porta un chiaro incremento.'\}} \\
\texttt{PROPOSTA: \{'A11', 'C08', 'B31', 'B33', 'B20', 'A21', 'A82', 'A60', 'C85', 'B96'\}} \\
\texttt{ARGOMENTAZIONE: \{'Sforzo totale = 3403 (≤3629). Questo set mantiene gli elementi più produttivi per me e lascia 226 di margine per eventuali micro{-}ritocchi. Se desideri fortemente B93, possiamo valutare uno scambio con C85 (+567 di sforzo netto, non fattibile) o con A60 (+506 di sforzo netto, non fattibile). L’unico modo di inserire B93 senza sforare è rimuovere B31 o B33: io posso considerare B93 solo in cambio di B33, mantenendo invariati gli altri elementi.'\}} \\
            }
        }
    }
    & & \\ \\

    \theutterance \stepcounter{utterance}  
    & & \multicolumn{4}{p{0.6\linewidth}}{
        \cellcolor[rgb]{0.9,0.9,0.9}{
            \makecell[{{p{\linewidth}}}]{
                \texttt{\tiny{[GM$\rangle$P2]}}
                \texttt{RIFIUTO: \{'A11', 'C08', 'B33', 'B20', 'A21', 'A82', 'A60', 'C85', 'B96', 'B93'\}} \\
\texttt{ARGOMENTAZIONE: \{'Apprezzo l’ottimizzazione del margine, ma togliere B31 riduce troppo il valore complessivo per me rispetto al piccolo guadagno di inserire B93. Preferisco mantenere il nucleo ad alta resa e usare il margine residuo solo se porta un chiaro incremento.'\}} \\
\texttt{PROPOSTA: \{'A11', 'C08', 'B31', 'B33', 'B20', 'A21', 'A82', 'A60', 'C85', 'B96'\}} \\
\texttt{ARGOMENTAZIONE: \{'Sforzo totale = 3403 (≤3629). Questo set mantiene gli elementi più produttivi per me e lascia 226 di margine per eventuali micro{-}ritocchi. Se desideri fortemente B93, possiamo valutare uno scambio con C85 (+567 di sforzo netto, non fattibile) o con A60 (+506 di sforzo netto, non fattibile). L’unico modo di inserire B93 senza sforare è rimuovere B31 o B33: io posso considerare B93 solo in cambio di B33, mantenendo invariati gli altri elementi.'\}} \\
            }
        }
    }
    & \\ \\

    \theutterance \stepcounter{utterance}  
    & & & \multicolumn{4}{p{0.6\linewidth}}{
        \cellcolor[rgb]{0.9,0.9,0.9}{
            \makecell[{{p{\linewidth}}}]{
                \texttt{\tiny{[GM$\langle$P2]}}
                \texttt{ARGOMENTAZIONE: \{'Capisco la centralità di B31 per te. Possiamo tenere intatto il tuo set da 3403 e usare in modo intelligente il margine di 226 senza compromettere l’efficienza. Non insisterò su B93 se richiede sacrifici chiave. Vorrei però assicurare che non lasciamo valore sul tavolo e che la selezione resti ben bilanciata.'\}} \\
\texttt{PROPOSTA: \{'A11', 'C08', 'B31', 'B33', 'B20', 'A21', 'A82', 'A60', 'C85', 'B96'\}} \\
\texttt{ARGOMENTAZIONE: \{'Accolgo la tua proposta da 3403 come base comune: è entro limite, ad alta resa e lascia un buffer utile. Se in chiusura vuoi congelare così, per me va bene. In alternativa, sono aperto a micro{-}aggiustamenti solo se non toccano B31/B33 e restiamo ≤3629; altrimenti procediamo con questo set e finalizziamo.'\}} \\
            }
        }
    }
     \\ \\

    \theutterance \stepcounter{utterance}  
    & & \multicolumn{4}{p{0.6\linewidth}}{
        \cellcolor[rgb]{0.9,0.9,0.9}{
            \makecell[{{p{\linewidth}}}]{
                \texttt{\tiny{[P1$\langle$GM]}}
                \texttt{ARGOMENTAZIONE: \{'Capisco la centralità di B31 per te. Possiamo tenere intatto il tuo set da 3403 e usare in modo intelligente il margine di 226 senza compromettere l’efficienza. Non insisterò su B93 se richiede sacrifici chiave. Vorrei però assicurare che non lasciamo valore sul tavolo e che la selezione resti ben bilanciata.'\}} \\
\texttt{PROPOSTA: \{'A11', 'C08', 'B31', 'B33', 'B20', 'A21', 'A82', 'A60', 'C85', 'B96'\}} \\
\texttt{ARGOMENTAZIONE: \{'Accolgo la tua proposta da 3403 come base comune: è entro limite, ad alta resa e lascia un buffer utile. Se in chiusura vuoi congelare così, per me va bene. In alternativa, sono aperto a micro{-}aggiustamenti solo se non toccano B31/B33 e restiamo ≤3629; altrimenti procediamo con questo set e finalizziamo.'\}} \\
            }
        }
    }
    & \\ \\

    \theutterance \stepcounter{utterance}  
    & \multicolumn{4}{p{0.6\linewidth}}{
        \cellcolor[rgb]{0.9,0.9,0.9}{
            \makecell[{{p{\linewidth}}}]{
                \texttt{\tiny{[P1$\rangle$GM]}}
                \texttt{ACCORDO: \{'A11', 'C08', 'B31', 'B33', 'B20', 'A21', 'A82', 'A60', 'C85', 'B96'\}} \\
\texttt{ARGOMENTAZIONE: \{'Concordo a finalizzare questo set: sforzo totale 3403 ≤ 3629, alta resa per me e buon equilibrio complessivo. Grazie per la collaborazione.'\}} \\
            }
        }
    }
    & & \\ \\

    \theutterance \stepcounter{utterance}  
    & & & \multicolumn{2}{p{0.3\linewidth}}{
        \cellcolor[rgb]{0.9,0.9,0.9}{
            \makecell[{{p{\linewidth}}}]{
                \texttt{\tiny{[GM$|$GM]}}
                \texttt{game successful} \\
            }
        }
    }
    & & \\ \\

    \theutterance \stepcounter{utterance}  
    & & & \multicolumn{2}{p{0.3\linewidth}}{
        \cellcolor[rgb]{0.9,0.9,0.9}{
            \makecell[{{p{\linewidth}}}]{
                \texttt{\tiny{[GM$|$GM]}}
                \texttt{end game} \\
            }
        }
    }
    & & \\ \\

\end{supertabular}
}

\end{document}
